\section{Evaluation}
\label{sec:evaluation}

This chapter describes the evaluation of the \texttt{visionOS}-based prototype, focusing on system performance, interaction quality, and perceived pose feedback in realistic usage scenarios. Tests were carried out both in the Xcode simulator and on the physical Apple Vision Pro, with the full backend pipeline active (AirPlay capture, main API, and FoundationPose-based pose estimation).  

Quantitative benchmarking is still ongoing, but the system has been instrumented with timestamps and logging to support later analysis. The following sections describe the test setups, qualitative results, and implications for future work. Additional plots and screenshots are provided in Appendix~\ref{appendix:evaluationfigures}.

\subsection{Metrics and Test Setups}
\label{subsec:evaluation-metrics}

Three evaluation dimensions were defined, summarized in Table~\ref{tab:evaluation-dimensions}.

\begin{table}[H]
    \centering
    \caption{Evaluation dimensions}
    \label{tab:evaluation-dimensions}
    \begin{tabular}{p{0.25\linewidth}p{0.63\linewidth}}
        \toprule
        \textbf{Dimension} & \textbf{Description} \\
        \midrule
        System performance &
        End-to-end responsiveness, API round-trip time, and visual smoothness of pose updates. \\
        User interaction and usability &
        Ease of ROI definition, stability of gaze-based selection, window layout, and perceived workload. \\
        Pose feedback quality &
        Alignment of the overlay with the physical object and temporal stability under motion and occlusion. \\
        \bottomrule
    \end{tabular}
\end{table}

These dimensions were probed in three controlled test setups; Table~\ref{tab:testmatrix} summarizes the scenarios.

\begin{table}[H]
    \centering
    \caption[Test matrix for prototype evaluation]{Test setups used to probe system performance, interaction quality, and pose feedback.}
    \label{tab:testmatrix}
    \begin{tabular}{@{}p{0.08\linewidth}p{0.26\linewidth}p{0.28\linewidth}p{0.28\linewidth}@{}}
        \toprule
        \textbf{Setup} & \textbf{Platform / environment} & \textbf{Objects and background} & \textbf{Primary focus} \\
        \midrule
        A &
        Xcode \texttt{visionOS} simulator, synthetic view &
        Static cube-like proxy object in a simple scene &
        API integration, polling stability, UI wiring \\
        B &
        Real AVP, indoor lab environment &
        Mac~Mini box as cube proxy; static placement, low clutter &
        Baseline latency, ROI control, static pose alignment \\
        C &
        Real AVP, indoor lab, mixed background &
        Banana and spanner models, moderate clutter, user and object motion &
        Pose stability, motion artefacts, robustness under stress \\
        \bottomrule
    \end{tabular}
\end{table}

In all setups, the same backend pipeline was used: AirPlay mirroring to a host machine, frame capture and mask extraction on the main API, depth estimation (monocular or external depth, depending on configuration), and outsourced 6D pose estimation via the FoundationPose-based service. The VisionOS client continuously streamed headset pose to the backend and requested corrected object poses via \texttt{/avp\_pose}.  

Representative screenshots of simulator tests, real-device overlays for the Mac~Mini box, banana, and spanner, as well as example masks are provided in Appendix~\ref{appendix:evaluationfigures}.

\subsection{System Performance}

System performance was assessed qualitatively in terms of responsiveness and visual smoothness of pose updates.

\subsubsection{Simulator Tests (Setup A)}

In the Xcode simulator, neither network latency nor AirPlay mirroring are part of the loop. End-to-end response from pose request to overlay update is therefore dominated by local JSON encoding/decoding and rendering. Under these conditions:

\begin{itemize}
    \item pose updates appeared immediate to testers,
    \item the polling loop could run at high frequencies without affecting responsiveness,
    \item API contract and error handling (e.g.\ malformed responses) could be verified without external bottlenecks.
\end{itemize}

Setup~A was thus mainly used as a sanity check for the client–backend interface and the robustness of the polling loop under repeated start/stop cycles.

\subsubsection{Real-Device Tests (Setups B and C)}

On real hardware, end-to-end behavior is dominated by the components listed in Table~\ref{tab:latency-drivers-eval}.

\begin{table}[H]
    \centering
    \caption{Dominant latency drivers in real-device tests}
    \label{tab:latency-drivers-eval}
    \begin{tabular}{p{0.28\linewidth}p{0.60\linewidth}}
        \toprule
        \textbf{Stage} & \textbf{Impact} \\
        \midrule
        AirPlay mirroring and capture &
        Adds encoding/decoding and screen-grab latency; occasional stutter and tiling under network load. \\
        CV pipeline (depth + mask) &
        Processing on backend (MDE or external depth, color-based mask reconstruction). \\
        FoundationPose inference &
        GPU-bound 6D pose estimation; typically the single largest contributor per cycle. \\
        Network AVP \texorpdfstring{$\leftrightarrow$}{<->} backend &
        Round-trip delays for \texttt{/head\_pose} and \texttt{/avp\_pose} requests. \\
        \bottomrule
    \end{tabular}
\end{table}

In Setup~B (static Mac~Mini box), testers described pose updates as “slightly delayed but usable” for inspection and slow user motion. Perceived latency matched the rough breakdown from Section~\ref{subsec:cv-pipeline-design}, i.e.\ a few hundred milliseconds dominated by pose inference and AirPlay overhead.  

In Setup~C (banana and spanner with user motion), the same processing latency combined with higher motion led to visible temporal lag between object motion and overlay updates. AirPlay occasionally produced stutter and tiling artefacts under network load, which resulted in transient pose outliers. Debug window logs confirmed corresponding spikes in \texttt{/avp\_pose} round-trip times during these episodes.

\begin{figure}[H]
    \centering
    \fbox{\parbox[c][4cm][c]{10cm}{\centering Placeholder for latency/performance chart (simulator vs.\ real AVP, Setups B/C)}}
    \caption[Prototype performance characteristics]{Qualitative performance characteristics in simulator and real-device setups: perceived update rate and API round-trip variability.}
    \label{fig:performanceplot}
\end{figure}

Overall, the architecture was sufficient for interactive inspection and low-speed manipulation, but not yet suitable for tightly coupled closed-loop control.

\subsection{User Interaction and Usability}

Interaction quality was evaluated informally with internal testers performing repeated ROI definition and model selection tasks in all setups. No formal user study was conducted, but structured feedback was collected after each session.

\subsubsection{ROI Definition and Window Layout}

Across setups, testers were consistently able to:

\begin{itemize}
    \item position main, ROI, and debug windows in comfortable 3D arrangements,
    \item select the correct CAD proxy model (cube, banana, spanner) from the main window,
    \item adjust ROI radius and color until the backend mask matched the intended region.
\end{itemize}

The gaze-anchored ROI center was stable under slow head motion. Misalignments occurred mainly when:

\begin{itemize}
    \item users moved their head quickly, or
    \item the gaze target was close to the edge of the mirrored AirPlay frame.
\end{itemize}

In these cases, the pixel coordinates inferred for the ROI center could drift, leading to masks that partially missed the object. Testers noted that having ROI and debug windows visible simultaneously made it easier to tune color thresholds and radius, because changes in the UI could be immediately correlated with mask images and pose success/failure messages in the logs.

\subsubsection{Perceived Usability}

Table~\ref{tab:ux-findings-eval} summarizes the main usability observations reported by testers.

\begin{table}[H]
    \centering
    \caption{Qualitative usability observations}
    \label{tab:ux-findings-eval}
    \begin{tabular}{p{0.30\linewidth}p{0.58\linewidth}}
        \toprule
        \textbf{Aspect} & \textbf{Observation} \\
        \midrule
        Window separation &
        Clear separation between configuration (main window), ROI tuning (ROI window), and diagnostics (debug window) was perceived as helpful. \\
        ROI feedback timing &
        ROI changes only visible after the next polling cycle; users requested clearer indication when a new pose request is in flight. \\
        Control flow &
        A dedicated “freeze/snapshot” button was requested to decouple ROI definition from continuous streaming, especially in Setup~C. \\
        Learnability &
        Interface considered understandable after a short familiarization phase; main concepts (model, ROI, polling) were easy to grasp. \\
        \bottomrule
    \end{tabular}
\end{table}

\begin{figure}[H]
    \centering
    \fbox{\parbox[c][5cm][c]{10cm}{\centering Placeholder for user experience ratings (bar graph across Setups A/B/C)}}
    \caption[User ratings on usability and intuitiveness]{Indicative internal user ratings on usability and intuitiveness across test setups.}
    \label{fig:uxratings}
\end{figure}

These findings point to targeted improvements in ROI feedback and interaction timing, without changes to the overall architecture.

\subsection{Pose Estimation Feedback Quality}

Pose feedback quality was judged by visual inspection of overlay alignment and temporal stability across setups.

\subsubsection{Static Conditions (Setups A and B)}

In static scenes (simulator, and real AVP with the Mac~Mini box), the arrow overlay was generally well aligned with the visible object edges when:

\begin{itemize}
    \item the ROI correctly enclosed the target, and
    \item the AirPlay frame was free of artefacts.
\end{itemize}

Under these conditions:

\begin{itemize}
    \item successive pose updates showed only minor jitter,
    \item the arrow appeared anchored to the box during slow head translations and rotations,
    \item orientation changes were plausible when the ROI or viewpoint was adjusted.
\end{itemize}

Residual misalignment could typically be traced to one of the causes in Table~\ref{tab:misalignment-causes-eval}.

\begin{table}[H]
    \centering
    \caption{Observed causes of residual pose misalignment}
    \label{tab:misalignment-causes-eval}
    \begin{tabular}{p{0.30\linewidth}p{0.58\linewidth}}
        \toprule
        \textbf{Cause} & \textbf{Effect} \\
        \midrule
        Head pose/frame mismatch &
        Small time offset between headset pose used for correction and frame used for pose estimation; results in slight spatial drift. \\
        Calibration inaccuracies &
        Imperfect intrinsics and ArUco-based calibration in the AirPlay setup; leads to systematic offset. \\
        Partial occlusions &
        Objects briefly occluding the Mac~Mini box; sometimes produces inconsistent pose estimates. \\
        \bottomrule
    \end{tabular}
\end{table}

\subsubsection{Motion and Stress Conditions (Setup C)}

With banana and spanner models in a cluttered environment and deliberate user motion, the system was exposed to:

\begin{itemize}
    \item more complex object geometry and texture,
    \item stronger perspective changes as the user moved,
    \item intermittent occlusions by hands or other objects.
\end{itemize}

In these conditions:

\begin{itemize}
    \item pose updates showed small but noticeable jumps, especially when the ROI partially covered background,
    \item AirPlay artefacts occasionally triggered transient outliers in the pose,
    \item no heavy temporal filtering was applied online, so jitter was directly visible in the overlay.
\end{itemize}

Testers reported that, although the overall direction and position remained reasonable, these discontinuities reduced subjective confidence in the pose during fast motion.

\begin{figure}[H]
    \centering
    \fbox{\parbox[c][5cm][c]{10cm}{\centering Placeholder for overlay comparison (Mac~Mini vs.\ banana vs.\ spanner, static vs.\ motion)}}
    \caption[Pose overlay comparison in different motion conditions]{Pose overlay comparison for different objects and motion patterns.}
    \label{fig:poseoverlay}
\end{figure}

\subsection{Simulation-Based Occlusion and Filtering Experiments}

In parallel to live AVP tests, the occlusion-aware filtering methods from Appendix~\ref{appendix:willems} were evaluated in simulation. These experiments operate purely on:

\begin{itemize}
    \item the sequence of previously estimated object poses, and
    \item the streamed headset pose data,
\end{itemize}

without access to user inputs or raw images.

\subsubsection{Filter Models}

Three classes of filters were compared:

\begin{itemize}
    \item \textbf{Extended Kalman Filter (EKF)} – pose state $\mathbf{x}_k$ with nonlinear process
    \[
        \mathbf{x}_{k+1} = f(\mathbf{x}_k, \mathbf{u}_k) + \mathbf{w}_k,
        \qquad
        \mathbf{y}_k = h(\mathbf{x}_k) + \mathbf{v}_k,
    \]
    where $\mathbf{u}_k$ encodes headset motion, and $\mathbf{w}_k$, $\mathbf{v}_k$ are process and measurement noise. Linearization yields the standard Kalman gain update.
    \item \textbf{Particle Filter (PF)} – posterior $p(\mathbf{x}_k \mid \mathbf{y}_{1:k})$ represented by weighted samples $\{ \mathbf{x}_k^{(i)}, w_k^{(i)} \}$; particles are propagated by the same motion model and reweighted according to measurement likelihood.
    \item \textbf{Gaussian Process (GP) with data-driven dynamics} – each pose component modelled as
    \[
        x(t) \sim \mathcal{GP}(m(t), k(t,t')),
    \]
    with temporal kernel $k(\cdot,\cdot)$; Willems’ lemma is used to construct a data-driven linear representation from recorded trajectories, which then feeds a GP residual model (see Appendix~\ref{appendix:willems}).
\end{itemize}

In all cases, the filters take only headset pose and raw object pose estimates as input; user gestures and task context do not enter the model.

\subsubsection{Qualitative Findings}

Simulation runs with synthetic occlusions and measurement dropouts showed:

\begin{itemize}
    \item EKF effectively reduces small, approximately Gaussian perturbations but is sensitive to model mismatch.
    \item PF is more robust to multi-modal uncertainty and abrupt changes, but computationally heavier.
    \item The GP-based approach, given sufficiently rich (persistently exciting) trajectories, provides smooth corrections and handles slow drift well, at the cost of a more expensive offline training phase.
\end{itemize}

These findings support the idea of inserting a filter stage between FoundationPose and \texttt{/avp\_pose} in the backend. The VisionOS client would remain unchanged, as the interface continues to be a 4$\times$4 transformation matrix.

\subsection{Discussion}

The current evaluation shows that the prototype supports interactive, user-driven AR inspection and low-speed interaction, despite:

\begin{itemize}
    \item ADP-level access restrictions on Apple Vision Pro,
    \item reliance on AirPlay for frame acquisition,
    \item external, GPU-bound pose inference.
\end{itemize}

The main limitations observed are:

\begin{itemize}
    \item end-to-end latency dominated by AirPlay and external pose estimation,
    \item pose jitter and occasional outliers when the stream is degraded or the ROI includes clutter,
    \item sensitivity of ROI placement to fast head motion and gaze inaccuracies.
\end{itemize}

At the same time, the modular design makes these issues addressable:

\begin{itemize}
    \item backend optimizations or moving depth/pose closer to the sensor (e.g.\ on-device inference) would directly reduce latency,
    \item integrating EKF/PF/GP-based filters in the backend could stabilize pose trajectories without touching the VisionOS client,
    \item improving ROI feedback and timing on the client side would make interaction more robust.
\end{itemize}

Future evaluations with a larger and more diverse user group, longer sessions, and more complex tasks (e.g.\ actual robot teaching workflows) will be required to fully characterize comfort, learnability, and long-term robustness. The current results demonstrate that the end-to-end design is functional and that the main technical risks—VisionOS integration, backend connectivity, and pose overlay—have been resolved in a working prototype that can now serve as a platform for further refinement.
